\chapter*{Abstract}
\addcontentsline{toc}{chapter}{Abstract}

Mutual information measures statistical dependence without restricting that
dependence to a linear or monotone form. Standard significance procedures for
categorical mutual information primarily test whether two variables are
independent within one population. This thesis studies a different problem:
whether two independent populations have equal mutual information,
\(H_0:I(P)=I(Q)\), while allowing their complete joint distributions to differ.

The usual analytic approach forms a bias-corrected difference of plug-in
mutual information estimates, divides it by an influence-function standard
error, and compares the resulting statistic with a standard normal
distribution. This normal Wald reference does not represent finite-sample
uncertainty in the estimated variance. The thesis develops an expanded
Welch--Satterthwaite correction that retains the same effect estimate and
standard error but derives mutual-information-specific component degrees of
freedom from the sampling variability of the complete variance functional.
The resulting procedure is deterministic and requires \(O(rc)\) operations
for an \(r\times c\) table.

The method was evaluated using 60 fixed pairs of positive-support categorical
populations with exactly equal mutual information. The experiment covered six
table shapes, five sampling regimes, sample sizes from 50 to 1,000, and 10,000
independent table-pair replicates per population pair. At nominal
\(\alpha=0.05\), mean absolute calibration error across all scenarios was
0.01342 for normal Wald, 0.01194 for simple Welch--Satterthwaite, and 0.00825
for expanded Welch--Satterthwaite. The expanded method left the well-sampled
control essentially unchanged and reduced error in the skewed, rare-cell, and
widespread-sparsity regimes, although it did not restore exact calibration in
the most difficult cases. Its average power loss relative to normal Wald was
approximately one percentage point, and median runtime remained below
0.15 milliseconds for the tested tables.

These results support expanded Welch--Satterthwaite inference as a low-cost
finite-sample correction for comparing positive mutual information values.
The method is not a universal sparse-table solution and does not apply at the
nonregular independence boundary, where the first-order mutual-information
variance vanishes.
